\SourceLeaf{173}\bigskip
\Needspace{10\baselineskip}
\centerline{[3.]}

\smallskip

\centerline{\textls{Ideen.}}

\smallskip

Die schärfste Bestimmung der Declinationen der Sterne würde in der Nähe des Äquators durch gute Repetitionstheodolithen gemacht werden können, indem ihre Azimuthalunterschiede mit terrestrischen Objecten zur Zeit der grössten östlichen und westlichen Digressionen beobachtet würden. \hfill\mbox{1813 April 21.}

\smallskip

In der Theorie der Parallellinien sind wir jetzt noch nicht weiter als Euklid war. Diess ist die partie honteuse der Mathematik, die früh oder spät eine ganz andere Gestalt bekommen muss. \hfill\mbox{1813 April 27.}

\smallskip

Eben so wie eine Classe von Wahrheiten der höhern Arithmetik, auf reelle Zahlen beschränkt, in inniger Verbindung mit den Kreisfunctionen steht, eben so werden sich die allgemeinern Wahrheiten der H.\,A., auf imaginäre Grössen ausgedehnt, in den Lemniscatischen Functionen gleichsam spiegeln. Diese Ideen öffnen die Pforte zu einem höchst reichhaltigen Felde der Analyse.

\smallskip

Das Krystallprisma an den Fernröhren vergrössert deren Länge etwa um $\frac{1}{3}$
 der Seite des Prisma, stört aber etwas die vollkommene Farbenlosigkeit.

\bigskip
\hrule
\bigskip

\centerline{BEMERKUNG.}

\smallskip

\textsc{Schumacher} hat sich während des Winters 1808 bis 1809 in Göttingen aufgehalten. Über \textsc{Gauss}’ Gespräche mit ihm hat er Aufzeichnungen in einem »GAUSSIANA« betitelten Hefte gemacht. Aus diesem Hefte ist die Notiz [2.] entnommen, während die Notiz [3.] auf einem einzelnen Zettel verzeichnet ist.
\begin{flushright}\textsc{Stäckel}.\end{flushright}

\bigskip
\hrule
